\section*{5 Full PRG via Iterated Restrictions}
So far, we have shown how to pseudorandomly assign values to a constant fraction of the inputs of any read-once PARITY $\circ$ AND formula using $O(\log n)$ truly random bits, preserving the expectation of the formula to within near-optimal error. In this section, to complete the proof of Theorem1, we show how to pseudorandomly assign values to all the inputs, i.e., we give a genuine PRG.
For convenience, we make the following definitions.
\begin{authorclaim}{Definition26}
Let $w>0$. A $w$-proper formula is a read-once PARITY $\circ$ AND formula of width at most $w$ and length most $2^{8w}$. We say that such a formula is short if its length is at most $2^{4w}$; otherwise, we say that the formula is long.
\end{authorclaim}
Our main goal is to fool $(C\log n)$-proper formulas, but along the way, we will obtain a PRG for $w$-proper formulas with seed length $O(w)$ and error $\exp(-\widetilde\Omega(w))$, even for $w$ substantially smaller than $\log n$.
\subsection*{5.1 Restrictions for Proper Formulas}
Recall that Lemma9 provides a pseudorandom restriction that uses only $O(\log n)$ truly random bits. We now generalize this fact in two respects. First, in the case of $w$-proper formulas ($\log\log n\le w\le C\log n$), we improve the seed length to $O(w)$. Second, in the case of short $w$-proper formulas, we argue that the restriction simplifies the formula, in the sense that it transforms it into a $(w/2)$-proper formula.
\begin{authorclaim}{Lemma27}
For every $w,n\in\mathbb N$ with $w\le C\log n$, there is a distribution $X$ over $\{0,1,?\}^n$ with the following properties.
\begin{enumerate}
\item (Seed length) There is an explicit algorithm to sample from $X$ using just $O(w+\log\log n)$ truly random bits.
\item (Expectation preservation) If $f$ is a $w$-proper formula, then $X$ preserves the expectation of $f$ with error $\exp(-\Omega(w/\log w))$.
\item (Simplification) If $f$ is a short $w$-proper formula, then $\Pr[f|_X\text{ is a }(w/2)\text{-proper formula}]\ge1-2^{-w}$.
\end{enumerate}
\end{authorclaim}
\begin{proof}
Let $n'=2^{8w}\cdot w$. Let $Y$ be an $n'$-wise $\delta^3$-biased distribution where $\delta=(n')^{-12\bar C}$, and let $Z$ be $\gamma$-almost $k$-wise independent with marginals $1-2^{-C}$, where $k=6\log n'$ and $\gamma=(n')^{-9}$. Our restriction is $X=\operatorname{Res}(Y,Z)^{\circ2^{C+4}}$.
By standard constructions [25,3] and Claim7, $X$ can be explicitly sampled using $O(w+\log\log n)$ truly random bits.
Now, to prove expectation preservation, let $f$ be a $w$-proper formula. By $w$-properness, there is some set of indices $I\subseteq[n]$, $|I|\le n'$, such that $f(x)$ only depends on $x|_I$. Let $g:\{0,1\}^{|I|}\to\{0,1\}$ be the $w$-proper formula such that $f(x)=g(x|_I)$. Since $Y|_I$ is $\delta^3$-biased and $Z|_I$ is $\gamma$-almost $k$-wise independent with marginals $1-2^{-C}$, Lemma9 implies that $\operatorname{Res}(Y|_I,Z|_I)$ preserves the expectation of $g$ with error $\exp(-\Omega(\log n'/\log\log n'))$, which is $\exp(-\Omega(w/\log w))$. It follows that $\operatorname{Res}(Y,Z)$ preserves the expectation of $f$ with the same error. The error of $X$ is only larger by a constant factor $2^{C+4}$, because any restriction of a $w$-proper formula is trivially another $w$-proper formula.
Finally, to prove simplification, let $f$ be a short $w$-proper formula, and let $f_i$ be a term. Since $k>w/2$, by Claim8, the probability that more than $w/2$ variables from $f_i$ are assigned ? by $X$ is bounded by
\begin{align*}
\binom w{w/2}\cdot((1-2^{-C})^{2^{C+4}\cdot w/2}+2^{C+4}\gamma)
&\le2^w\cdot(e^{-2^{-C}\cdot2^{C+3}w}+2^{C+4}\cdot(n')^{-9})\\
&<2^{-5w}.
\end{align*}
The number of terms in $f$ is at most $2^{4w}$, so by the union bound, except with probability $2^{-w}$, $f|_X$ has maximum width at most $w/2$. Furthermore, restricting cannot increase the number of terms, so the number of terms is still bounded by $2^{4w}=2^{8(w/2)}$. Therefore, in this case, $f|_X$ is $(w/2)$-proper.
\end{proof}
\subsection*{5.2 Full PRGs for Long Proper Formulas [24]}
The simplification clause of Lemma27 only applies if $f$ is short. If $f$ is long, we will therefore need a different approach. We will take a similar approach as Meka, Reingold, and Tal [24]. A full PRG for long $w$-proper formulas follows readily from their work.
\begin{authorclaim}{Lemma28 ([24])}
For every $w,n\in\mathbb N$, there is an explicit $(2^{-w})$-PRG for long $w$-proper formulas with seed length $O(w+\log\log n)$.
\end{authorclaim}
\begin{proof}[Proof sketch]
In short, the PRG is one of the PRGs by Meka et al. [24, full version, Lemma6.2], except we replace every $\delta$-biased distribution with a $(\cdot)$-wise $\delta$-biased distribution to optimize the seed length.
In more detail, let $n'=2^{8w}\cdot w$. Sample $v\in\{0,1\}^{wn}$ from an $(n'w)$-wise $(c_{\mathrm{MRT}}^{-w})$-biased distribution, where $c_{\mathrm{MRT}}$ is a suitable constant. Think of $v$ as $n$ blocks of $w$ bits. Define a set $I\subseteq[n]$ as follows: include $i$ in $I$ if and only if the $i$-th block of $v$ is $1^w$.
Sample $x^{(0)},x^{(1)},\ldots,x^{(16)}\in\{0,1\}^n$ independently from an $(n')$-wise $(c_{\mathrm{MRT}}^{-w})$-biased distribution. The PRG outputs the string $x$ defined by $x_i=x_i^{(0)}$ if $i\notin I$, and $x_i=\bigoplus_{j=1}^{16}x_i^{(j)}$ if $i\in I$.
By standard constructions [25,3], the seed length of this PRG is $O(\log n'+w+\log\log n)=O(w+\log\log n)$.
As for correctness, let $f$ be a long $w$-proper formula. Let $J\subseteq[n]$ be the set of indices of variables that $f$ reads, so there is some long $w$-proper formula $g$ on $|J|$ input bits such that $f(x)=g(x|_J)$. Let $X$ be the distribution output by the PRG. Since $|J|\le n'$, the distribution $X|_J$ is exactly the pseudorandom distribution designed by Meka et al. [24, full version, Lemma6.2]. Furthermore, since $f$ is long, $|J|>2^{4w}$. It follows that $g$ is in the class of functions fooled by Meka et al.'s pseudorandom distribution: $g$ is an XOR of $m$ non-constant Boolean functions on disjoint variables, where each function is on at most $w$ variables, with $16^w<m\le16^{2w}$ and $\log\log(|J|/2^w)\ll w\le\log|J|$. Therefore, $X|_J$ fools $g$ with error $2^{-w}$, and hence $X$ fools $f$ with error $2^{-w}$.
\end{proof}
\subsection*{5.3 Full PRGs for Width-$O(\log n)$ Formulas}
For short proper $w$-formulas, to get a full PRG, we will iterate the restriction of Lemma27 several times, assigning values to more and more variables. Eventually, we'll stop this recursive process and use a different PRG. Specifically, for the "base case," we'll use a PRG by Lee [22] with minor modifications:
\begin{authorclaim}{Lemma29 ([22])}
For every $w,n\in\mathbb N$ and every $\epsilon>0$, there is an explicit $\epsilon$-PRG for $w$-proper formulas with seed length
\[
O((w+\log(1/\epsilon))\cdot(\log w+\log\log(1/\epsilon))^2)
+\operatorname{poly}(\log\log(n/\epsilon)).
\]
\end{authorclaim}
\begin{proof}[Proof sketch]
In short, the PRG is one of the PRGs by Lee [22, Theorem6], except we replace every $\delta$-biased distribution with a $(\cdot)$-wise $\delta$-biased distribution to optimize the seed length, just like the proofs of Lemma27 and Lemma28.
To give a little more detail, let $n'=2^{8w}\cdot w$; a $w$-proper formula only reads $n'$ variables. Lee's PRG [22, Theorem6] is designed to fool arbitrary-order combinatorial checkerboards, i.e., parities of functions on disjoint variable sets of size at most $w$. This class includes $w$-proper formulas as a special case. Lee's original PRG has seed length $O((w+\log(n/\epsilon))\cdot(\log w+\log\log(n/\epsilon))^2)$.
After making suitable replacements, one can show that the seed length is reduced to
\[
O((w+\log(n'/\epsilon))\cdot(\log w+\log\log(n'/\epsilon))^2)
+\operatorname{poly}(\log\log(n/\epsilon)).
\]
(We omit the full proof, since it repeats much of Lee's analysis [22].) Plugging in the value of $n'$, we get the claimed seed length.
\end{proof}
We now give our full PRG for general formulas of width at most $C\log n$. The PRG follows a similar approach to one of the PRGs by Meka et al. [24, full version, Algorithm3]: iteratively apply the restriction of Lemma27, but at each step, XOR with the PRG of Lemma28 in case the formula is long.
\begin{authorclaim}{Lemma30}
For every $n\in\mathbb N$, there is an explicit PRG for read-once PARITY $\circ$ AND formulas of width at most $C\log n$ with seed length $O(\log n)$ and error $2^{-\Omega(\log n/(\log\log n)^3)}$.
\end{authorclaim}
\begin{proof}
Define $w_0=\log n/(\log\log n)^2$. We recursively define a PRG $G_w$ for $w$-proper formulas, $w_0\le w\le C\log n$, as follows.
\begin{itemize}
\item (Base case) If $w\le2w_0$, then $G_w$ is the $(2^{-w_0})$-PRG of Lemma29 based on Lee's work [22].
\item (Recursive case) If $w>2w_0$, sample $X\in\{0,1,?\}^n$ from the distribution guaranteed by Lemma27 based on the work in Section4. Sample $Y\in\{0,1\}^n$ using the PRG of Lemma28 based on Meka et al.'s work [24]. Recursively sample $G_{\lceil w/2\rceil}$, and set $G_w=Y\oplus(X\circ G_{\lceil w/2\rceil})$.
\end{itemize}
For the analysis, observe first that in the base case $w\le2w_0$, $G_w$ fools $w$-proper formulas with error $2^{-w_0}$. Now, for the inductive step, consider some $w>2w_0$. Assume $G_{\lceil w/2\rceil}$ fools $\lceil w/2\rceil$-proper formulas with error $\epsilon_{\lceil w/2\rceil}$; we will show that $G_w$ fools $w$-proper formulas with error $\epsilon_w$, where $\epsilon_w=\epsilon_{\lceil w/2\rceil}+2^{-\Omega(w/\log w)}$.
Let $f$ be a $w$-proper formula, and for brevity, let $G=G_{\lceil w/2\rceil}$. For the first case, suppose $f$ is long. Any shift of $f$ is also a long $w$-proper formula, so
\begin{align*}
|\mathbb E[f(G_w)]-\mathbb E[f]|
&=\left|\mathbb E_{X,G}[\mathbb E_Y[f(Y\oplus(X\circ G))]]-\mathbb E[f]\right|\\
&\le\mathbb E_{X,G}[|\mathbb E_Y[f(Y\oplus(X\circ G))]-\mathbb E[f]|]\\
&=\mathbb E_{X,G}[|\mathbb E_Y[f(Y\oplus(X\circ G))]-\mathbb E_U[f(U\oplus(X\circ G))]|]\\
&\le2^{-w}.
\end{align*}
For the second case, suppose $f$ is short. For each $y\in\{0,1\}^n$, define $f_y(x)=f(y\oplus x)$, another short $w$-proper formula. Fix $y\sim Y$, and let $\mathcal E$ be the event that $f_y|_X$ is $(w/2)$-proper, so whether $\mathcal E$ occurs depends only on $X$. Then
\begin{align*}
|\mathbb E[(f_y|_X)(G)]-\mathbb E[f]|
&\le|\mathbb E_X[\mathbb E_G[(f_y|_X)(G)]\mid\mathcal E]-\mathbb E[f]|+\Pr[\neg\mathcal E]\\
&\le|\mathbb E_X[\mathbb E_U[(f_y|_X)(U)]\mid\mathcal E]-\mathbb E[f]|
+\epsilon_{\lceil w/2\rceil}+\Pr[\neg\mathcal E]\quad(\text{Induction})\\
&\le|\mathbb E_X[\mathbb E_U[(f_y|_X)(U)]]-\mathbb E[f]|
+\epsilon_{\lceil w/2\rceil}+2\Pr[\neg\mathcal E]\\
&\le2^{-\Omega(w/\log w)}+\epsilon_{\lceil w/2\rceil}+2\Pr[\neg\mathcal E]\quad(\text{Item2 of Lemma27})\\
&\le2^{-\Omega(w/\log w)}+\epsilon_{\lceil w/2\rceil}+2\cdot2^{-w}\quad(\text{Item3 of Lemma27}).
\end{align*}
Let $\epsilon_w$ be the final right-hand side, so indeed $\epsilon_w=\epsilon_{\lceil w/2\rceil}+\exp(-\Omega(w/\log w))$. Then
\[
|\mathbb E[f(G_w)]-\mathbb E[f]|
\le\mathbb E_Y[|\mathbb E_{X,G}[(f_Y|_X)(G)]-\mathbb E[f]|]\le\epsilon_w.
\]
Now, let us add up all these errors. Since $w\ge w_0$ always holds, we have $\epsilon_w\le\epsilon_{\lceil w/2\rceil}+\exp(-\Omega(w_0/\log w_0))$. Starting at $w=C\log n$, we only need to halve $w$ a total of $O(\log\log\log n)$ times to reach the base case $w\le2w_0$. Therefore, the total error of $G_{C\log n}$ is bounded by
\[
2^{-w_0}+2^{-\Omega(w_0/\log w_0)}\cdot O(\log\log\log n)
=2^{-\Omega(\log n/(\log\log n)^3)}.
\]
Finally, let us bound the seed length of $G_w$. In the base case $w\le2w_0$, by our choice of $w_0$, the seed length $s_w$ of $G_w$ is bounded by some value $s_{\mathrm{base}}\le O(\log n)$. In the recursive case $w>2w_0$, the seed length $s_w$ of $G_w$ is bounded by
$s_w=s_{\lceil w/2\rceil}+O(w+\log\log n)=s_{\lceil w/2\rceil}+O(w)$.
The point is that this is essentially a geometric series. More precisely, let $c_{\mathrm{seed}}$ be a constant such that $s_w\le s_{\lceil w/2\rceil}+c_{\mathrm{seed}}\cdot w$ for all $w>2w_0$. Then by induction, for all $w\ge w_0$, we have $s_w\le s_{\mathrm{base}}+3c_{\mathrm{seed}}w$, because
\begin{align*}
s_w&\le s_{\lceil w/2\rceil}+c_{\mathrm{seed}}w\\
&\le s_{\mathrm{base}}+3c_{\mathrm{seed}}\lceil w/2\rceil+c_{\mathrm{seed}}w\quad(\text{Induction})\\
&<s_{\mathrm{base}}+3c_{\mathrm{seed}}w.
\end{align*}
Therefore, we can take the desired PRG to be $G_{C\log n}$, because $s_{C\log n}\le O(\log n)$, and any read-once PARITY $\circ$ AND formula of width at most $C\log n$ is $(C\log n)$-proper.
\end{proof}
\subsection*{5.4 Arbitrary-Error PRGs for Width-$O(\log(n/\epsilon))$ Formulas}
At this point, the main work of proving Theorem1 is complete. We just need to address three minor issues: small $\epsilon$, large width, and formulas not of the form PARITY $\circ$ AND. We begin by addressing the case of small $\epsilon$. Recall that we wish to achieve seed length $O(\log n)+\widetilde O(\log(1/\epsilon))$ for an arbitrary error $\epsilon$. This follows readily by combining the PRG of Lemma30 with Lee's PRG (Lemma29).
\begin{authorclaim}{Lemma31}
For any $n\in\mathbb N$, $\epsilon>0$, there is an explict $\epsilon$-PRG for read-once PARITY $\circ$ AND formulas of width at most $(C/2)\log(n/\epsilon)$ with seed length $O(\log n+\log(1/\epsilon)\cdot(\log\log(1/\epsilon))^5)$.
\end{authorclaim}
\begin{proof}
Let $\epsilon_0$ be the error parameter in Lemma30, so $\epsilon_0=\exp(-\Omega(\log n/(\log\log n)^3))$. If $\epsilon\ge\epsilon_0$, the PRG of Lemma30 works, because $(C/2)\log(n/\epsilon)<C\log n$. If $\epsilon<\epsilon_0$, use Lee's PRG [22], i.e., the $\epsilon$-PRG of Lemma29 for $((C/2)\log(n/\epsilon))$-proper formulas, which has seed length
\[
O(\log(n/\epsilon)\cdot(\log\log(n/\epsilon))^2)
\le O(\log(1/\epsilon)\cdot(\log\log(1/\epsilon))^5).
\]
\end{proof}
(In the proof of Lemma31, we could just as well have used Lee's original PRG [22, Theorem6] instead of the slightly modified version given by Lemma29.)
\subsection*{5.5 PRGs for Any Width}
In this section, we eliminate the assumption that the maximum width is bounded.
\begin{authorclaim}{Lemma32}
For all $n\in\mathbb N$ and $\epsilon>0$, there is an explicit $\epsilon$-PRG for read-once PARITY $\circ$ AND formulas on $n$ input bits with seed length $O(\log n+\log(1/\epsilon)\cdot(\log\log(1/\epsilon))^5)$.
\end{authorclaim}
\begin{proof}
Sample $G$ from the $(\epsilon/3)$-PRG for formulas of width $(C/2)\log(3n/\epsilon)$ guaranteed by Lemma31. Sample $Y$ from an $(\epsilon/(6n))$-biased distribution. Our final PRG outputs $H\mathrel{\mathop:}=G\oplus Y$.
To prove that this works, let $f$ be a read-once PARITY $\circ$ AND formula. Write $f=f'\oplus f''$, where every term in $f'$ has width at most $(C/2)\log(3n/\epsilon)$ and every term in $f''$ has width greater than $(C/2)\log(3n/\epsilon)$.
Since any shift of a width-$w$ read-once PARITY $\circ$ AND formula is another width-$w$ read-once PARITY $\circ$ AND formula, $H$ fools $f'$ with error $\epsilon/3$. Meanwhile, since each term $f''_i$ of $f''$ is a conjunction of more than $(C/2)\log(3n/\epsilon)$ literals, $\mathbb E[f''_i]\le(\epsilon/(3n))^{C/2}<\epsilon/(6n)$.
Furthermore, the $L_1$ norm of any conjunction of literals is1, and $H$ is $(\epsilon/(6n))$-biased, so by Claim5, $\mathbb E[f''_i(H)]<\epsilon/(3n)$. Therefore, by the union bound, for either distribution $X\in\{H,U\}$, $\mathbb E[f''(X)]<\epsilon/3$.
This allows us to bound the error of the final PRG as follows:
\begin{align*}
|\mathbb E[f(H)]-\mathbb E[f]|
&\le|\mathbb E[f(H)]-\mathbb E[f'(H)]|+|\mathbb E[f'(H)]-\mathbb E[f']|+|\mathbb E[f']-\mathbb E[f]|\\
&\le\mathbb E[|f(H)-f'(H)|]+|\mathbb E[f'(H)]-\mathbb E[f']|+\mathbb E[|f'-f|]\\
&=\mathbb E[f''(H)]+|\mathbb E[f'(H)]-\mathbb E[f']|+\mathbb E[f'']\\
&<\epsilon/3+\epsilon/3+\epsilon/3=\epsilon.
\end{align*}
\end{proof}
\subsection*{5.6 Proof of Theorem1}
In this section, we finally complete the proof of Theorem1 by showing that fooling read-once PARITY $\circ$ AND formulas is sufficient for fooling read-once depth-2 $\mathrm{AC}^0[\oplus]$:
\begin{authorclaim}{Lemma33}
Let $X$ be a distribution over $\{0,1\}^n$, and let $\epsilon>0$. If $X$ fools all read-once PARITY $\circ$ AND formulas with error $\epsilon$, then $X$ fools all read-once depth-2 $\mathrm{AC}^0[\oplus]$ formulas with error $2\epsilon$.
\end{authorclaim}
\begin{proof}
Let $f$ be a read-once depth-2 $\mathrm{AC}^0[\oplus]$ formula.
For the first case, suppose the output gate of $f$ is $\oplus$. By merging the output gate with any $\oplus$ children and introducing trivial $\wedge$ gates with fan-in1 as necessary, we see that without loss of generality, every child of the output gate is either $\wedge$ or $\vee$. By de Morgan's laws, it follows that either $f$ or $\neg f$ can be computed by a read-once PARITY $\circ$ AND formula. Either way, this implies that $X$ $\epsilon$-fools $f$.
For the second case, suppose the output gate of $f$ is $\wedge$, say $f=\bigwedge_{i=1}^mf_i$. Using the Fourier expansion of the $m$-input AND function, we get
\[
f=\sum_{I\subseteq[m]}\frac{(-1)^{|I|}}{2^m}\prod_{i\in I}(-1)^{f_i}
=\sum_{I\subseteq[m]}\frac{(-1)^{|I|}}{2^m}\cdot\left(1-2\cdot\bigoplus_{i\in I}f_i\right).
\]
By our analysis for the first case, $X$ fools $\bigoplus_{i\in I}f_i$ with error $\epsilon$. Therefore, by the triangle inequality,
\begin{align*}
|\mathbb E[f(X)]-\mathbb E[f]|
&\le\left|\sum_{I\subseteq[m]}\frac{(-1)^{|I|}\cdot(-2)}{2^m}
\left(\mathbb E[(\bigoplus_{i\in I}f_i)(X)]-\mathbb E[\bigoplus_{i\in I}f_i]\right)\right|\\
&\le\sum_{I\subseteq[m]}\frac2{2^m}\cdot\epsilon=2\epsilon.
\end{align*}
For the final case, suppose the output gate of $f$ is $\vee$. By de Morgan's laws, $\neg f$ can be computed by a read-once depth-2 $\mathrm{AC}^0[\oplus]$ formula with output gate $\wedge$. By our analysis for the second case, $X$ fools $\neg f$ with error $2\epsilon$, hence $X$ fools $f$ with the same error.
\end{proof}
